%! AP Statistics — Unit 3, Lesson 3.5 — Activity KEY
\documentclass[10pt]{article}
\usepackage{apstats-article}
\usepackage{apstats-key}

\begin{document}

\pageheader{Unit 3, Lesson 3.5}{Group Activity KEY: Designing Experiments}
\namepartnerperiod

\vspace{0.1in}

\textbf{Scenario A:} Memory supplement (flawed design).
\begin{practicebox}
\small
\begin{enumerate}[leftmargin=1.5em, itemsep=5pt]
  \item \ans{30 volunteers}
  \item \ans{Supplement type (supplement vs.\ no supplement)}
  \item \ans{Take the supplement; take nothing (no treatment)}
  \item \ans{Memory test score}
  \item \ans{Random assignment is violated.} Why: \ans{Participants self-selected their treatment; motivated participants chose the supplement.}
  \item \ans{Motivation is a confounding variable: participants who want to take the supplement may already study more or sleep better, which improves memory regardless of the supplement.}
  \item \ans{Randomly assign the 30 participants to supplement or placebo group; use a placebo pill so neither group knows which they received (single-blind or double-blind).}
\end{enumerate}
\end{practicebox}

\vspace{0.08in}

\textbf{Scenario B:} Headache medication.
\begin{practicebox}
\small
\begin{enumerate}[leftmargin=1.5em, itemsep=5pt]
  \item \ans{CRD (completely randomized design)}
  \item \ans{Double-blind (patients and doctors don't know which brand)}
  \item \ans{No control group (both groups get a real medication).} \ans{For comparing two active treatments, a placebo is not always necessary, though it would help measure absolute effectiveness.}
  \item \ans{Yes — if age affects response to the medication, blocking by age group would remove that variability and give a more precise estimate of the treatment difference.}
\end{enumerate}
\end{practicebox}

\vspace{0.08in}

\begin{practicebox}
\small
\textbf{Part 3 (classical music):} Accept designs that include:
\begin{enumerate}[leftmargin=1.5em, itemsep=4pt]
  \item \ans{Music group (studies with classical music) vs.\ silent group (no music).}
  \item \ans{Randomly assign students (e.g., coin flip, random number generator).}
  \item \ans{Confounds: time of day (hold constant); topic studied (use same material for all); study environment (use same room).}
  \item \ans{Quiz score immediately after study session.}
  \item \ans{Full double-blind is not feasible (students know if they hear music). Single-blind possible if the person grading doesn't know which group each quiz came from.}
  \item Diagram: All 30 students $\xrightarrow{\text{random}}$ Music (15) / Silent (15) $\to$ quiz.
\end{enumerate}
\end{practicebox}

\end{document}
